The linear period data utility for the Iterative Model Solution (IMS) can be activated by specifying the LINEAR\_PERIODDATA option in the IMS input file.  If activated, \mf will read linear period data according to the following description.

The linear period data file contains keyword PERIOD blocks that override selected settings from the LINEAR block of the IMS input file on a stress-period basis.  A PERIOD block applies to the stress period given in its header and to all subsequent stress periods until a later PERIOD block changes the settings.  Stress periods without a PERIOD block continue to use the running settings.  Only the following LINEAR-block keywords may appear in a PERIOD block: INNER\_DVCLOSE, INNER\_RCLOSE, INNER\_MAXIMUM, PRECONDITIONER\_LEVELS, PRECONDITIONER\_DROP\_TOLERANCE, and RELAXATION\_FACTOR.  Any other keyword will cause \mf to terminate with an error.

Changing the inner closure tolerances or the inner maximum takes effect immediately at no additional cost.  Changing a preconditioner control (PRECONDITIONER\_LEVELS, PRECONDITIONER\_DROP\_TOLERANCE, or RELAXATION\_FACTOR) causes \mf to re-resolve the preconditioner type and, when the preconditioner storage requirements change, to reallocate and refactor the preconditioner before the next solve.

\vspace{5mm}
\subsection{Structure of Blocks}
\lstinputlisting[style=blockdefinition]{./mf6ivar/tex/utl-imslinearperiod-period.dat}

\vspace{5mm}
\subsection{Explanation of Variables}
\begin{description}
\input{./mf6ivar/tex/utl-imslinearperiod-desc.tex}
\end{description}

\vspace{5mm}
\subsection{Example Input File}
\lstinputlisting[style=inputfile]{./mf6ivar/examples/utl-imslinearperiod-example.dat}
